\documentclass[11pt]{article}
\usepackage{amsmath,amssymb,amsthm}
\usepackage[T1]{fontenc}
\usepackage{lmodern}
\usepackage[a4paper,margin=25mm,top=23mm,bottom=25mm]{geometry}
\usepackage[hidelinks]{hyperref}
\hypersetup{pdftitle={An unbounded additive gain for half-turn no-three-in-line sets},pdfauthor={Working research note}}
\setlength{\emergencystretch}{1.2em}
\newtheorem{theorem}{Theorem}\newtheorem{lemma}[theorem]{Lemma}\newtheorem{proposition}[theorem]{Proposition}
\newcommand{\F}{\mathbb F_p}\newcommand{\G}{\mathcal G}\newcommand{\D}{\mathcal D}\newcommand{\R}{\mathcal R}\newcommand{\Z}{\mathbb Z}\newcommand{\htf}{F_{\mathrm{ht}}}\newcommand{\norm}[1]{\lVert#1\rVert_\infty}\newcommand{\Prob}{\mathbb P}\newcommand{\ch}{\chi}\newcommand{\TT}{\mathcal T}
\title{An unbounded additive gain for half-turn no-three-in-line sets on a sequence of square grids}
\author{Working proof for independent review}\date{19 November 2026}
\begin{document}\maketitle
\begin{abstract}
On infinitely many grids of side $4p+1$, with $p\equiv1\pmod{840}$ prime, we obtain a half-turn invariant no-three-in-line set of size at least $6p+2+\exp(\log p/(300\log\log p))$. The old conic ring is not frozen: whole components may change orientation with no loss of cardinality. The main step bounds collisions between multiplicative cosets used by ordinary integer secants. Once the abscissa cosets of the blocking endpoints are all distinct, each Boolean variable occurs at most twice, and Hall's theorem clears the remaining secants. Stewart's large prime factor theorem supplies the required sparse subgroups. This is a subsequential additive gain, not a change in the leading coefficient $3/2$ or a guarantee at every side. The argument is substantially AI-generated and has not been independently refereed.
\end{abstract}

\section{Claim and strategy}
A \emph{cap} is a set of integer points with no three distinct points on an ordinary Euclidean line. Let $\htf(n)$ denote the largest size of a cap in $\{0,\ldots,n-1\}^2$ invariant under $(x,y)\mapsto(n-1-x,n-1-y)$.
\begin{theorem}\label{thm:main} There are infinitely many primes $p\equiv1\pmod{840}$ for which
\[
 \htf(4p+1)\ \ge\ 6p+2+\exp\!\left(\frac{\log p}{300\log\log p}\right).
\]
In particular, given any $m$, infinitely many such boards admit $6p+2+2m$ points: a ring of $6(p-1)$ points, eight specified boundary points, and $m$ additional opposite pairs. The conic level and the ring orientations may be chosen for the board and the pairs.
\end{theorem}
This is a net gain over $6p+2$, not an append-only claim for an arbitrarily prescribed set. The four-period ring, its component flips and the two-literal formulation of insertion constraints are due to \cite{ring}; the rooted product identity is from \cite{root}. The classical unrestricted density construction is \cite{hjsw}. We include the boundary check from the earlier eight-point construction \cite{boundary} on the subprogression used here.

\paragraph{Proof map.} Independently sample candidate opposite pairs at a single common random conic level. A fixed positive fraction in expectation satisfy two finite character tests. We bound the chance that distinct blocking endpoints use the same multiplicative coset, distinguishing ordinary lines that become parallel modulo $p$ from those that do not. We discard conflicting \emph{candidate pairs}, before committing to an insertion; we never thin the final ring. Hall's theorem then chooses one compatible orientation. The numerical value 300 is only a convenient slack constant in the last prime-selection step. Write $L_p=1+\log p$; logarithms are natural and all implied constants are absolute.

\section{Full rings with independent orientation bits}
Let $p\equiv1\pmod{840}$ be prime, $k\in\F^*$ a nonsquare, $H=\langle -1,3\rangle\subset\F^*$, $h=|H|$, $I=(p-1)/h$. Both generators are squares. The cosets $A$ and $k/A$ in $\F^*/H$ form distinct pairs: their quadratic characters differ. Independently in each pair give one coset physical colour 0 and the other 1. For $u\ne0$ let $e(u)$ be the colour of its coset. Throughout the board $[0,4p]^2$, the fibre over $(u,k/u)$ has points $(x_0+ip,y_0+jp)$ with least positive representatives $x_0,y_0$ and $0\le i,j\le3$. Take the sites outside $\{1,2\}^2$ in sheet indices and having physical parity $x+y\equiv e(u)\pmod 2$. Writing $ij$ for the sheet pair $(i,j)$, the two possible masks are
\[
 M_0=\{00,02,13,20,31,33\},\qquad M_1=\{01,03,10,23,30,32\}.
\]
Each is a six-point cap. Subtracting its points in pairs shows that the only primitive nonaxial directions in a mask, up to sign, are $(1,\pm1),(3,\pm1),(1,\pm3)$. A primitive physical line reduces to a genuine line modulo $p$ and meets $xy=k$ in at most two residues. An old triple would therefore contain a same-fibre pair. An axis meets only one residue of the conic. For a nonaxial repeated direction $(a,b)$ the other abscissa is $-ak/(bu)$ by the quadratic root product. This lies in the coset paired with $uH$; its colour is opposite, whereas $x+y$ has constant parity on the physical line. Thus every orientation gives a cap $O$ of size $6(p-1)$. The reflection in $C=(2p,2p)$ sends $u$ to $-u$, preserves parity and the outer sheets, so $O$ has the required half-turn. Every noncoarse physical row and column strictly inside $(p,3p)$ has exactly one old point: the two outer sheet choices in the other coordinate have opposite parities.

Add the eight anchors $B=p\{(0,1),(0,3),(1,0),(3,0),(4,1),(4,3),(1,4),(3,4)\}$. They form a cap by convexity and the two-per-edge count. From an anchor, a nonaxial primitive direction $(a,b)$ meets the carrier only at parameters satisfying $ab t^2=k\pmod p$. At height $\max(|a|,|b|)\le7$ every nonzero signed factor is a square modulo $p$ by quadratic reciprocity, so there are no such points. At height $D\ge8$ only the inward ray occurs, with $1\le t\le4p/D\le p/2$; it contains at most one of the two opposite roots modulo $p$ (and none if $p\mid ab$). A coarse axis misses the carrier. Lines between anchors have primitive height at most four. Thus $O\cup B$ is a cap for \emph{every} choice of the orientation bits.

\paragraph{Matching criterion.} For a new site $q$ and two potentially selected outer-sheet endpoints on one secant through $q$, forbid simultaneous selection of the endpoints. Each endpoint is selected at exactly one value of its bit. Hence this is a two-literal clause: make at least one endpoint absent. We omit axes because the selected central rows/columns have only one old site; we shall insist that different targets use different rows and columns. For pairs of opposite new sites there is no need to repeat the clauses at the reflected sites. Suppose the abscissa $H$-cosets attached to \emph{all endpoint occurrences} of the remaining clauses are pairwise distinct. A bit belongs to two cosets, so it has at most two occurrences total. Any $j$ clauses have $2j$ occurrences and therefore at least $j$ distinct neighbouring bits, even when both endpoints of a clause use one bit. By Hall's marriage theorem match each clause to a different bit and set that bit to disable one of its endpoints. This satisfies all clauses simultaneously. This is why we need to separate equal cosets, \emph{not} all inverse-paired cosets.

\section{Sampling the targets and excluding short patterns}
Choose one level $k\in\F$ uniformly. Independently of it, sample $N$ independent uniform residues $z_i=(X_i,Y_i)\in\F^2$. Index $i$ proposes the two actual sites $q_i=(p+[X_i]_p,p+[Y_i]_p)$ and $2C-q_i$, where $[\ ]_p\in[0,p-1]$. We call $q_i$ the representative and $z_i$ its root modulo $p$. Write $d_i=X_iY_i$ and let $\ch$ be the quadratic character, zero at zero. We now define a test $E_i$ for \emph{one} index and show $\Prob(E_i)\ge\gamma>0$ for all sufficiently large $p$, with absolute $\gamma$. We will use linearity of expectation, not an assumption that these tests at the shared level are independent.

First require $\ch(d_i)=+1$ and $\ch(k)=-1$. For each primitive $(a,b)$ with $a>0$, $ab\ne0$, $\max(|a|,|b|)\le8$ require
\begin{equation}\label{eq:shortslope}
 \ch(k-\rho_{i,a,b})=-1,\qquad \rho_{i,a,b}=-(aY_i-bX_i)^2/(4ab).
\end{equation}
The discriminant of $(X_i+at)(Y_i+bt)=k$ is $4ab(k-\rho_{i,a,b})$, and $4ab$ is a square. In these directions no carrier point lies on the line, even before a mask is chosen.

Consider two carrier endpoints at signed offsets $t,u$ from $q_i$ along a primitive direction. Divide $t,u$ by their common gcd; the larger of the two resulting absolute values is the \emph{ratio height}. To exclude ratio heights at most 8, require $\ch(k-d_i)=-1$ and
\begin{equation}\label{eq:shortratio}
 \ch(k-\beta d_i)=+1\quad(\beta\in\R),\qquad
 \R=\left\{\left(\frac{r-s}{r+s}\right)^2:\ 0<|r|,|s|\le8,\ r\ne\pm s\right\}.
\end{equation}
Indeed let the reduced parameters be $r,s$, and $(U,V)$ the first endpoint modulo $p$. The second is $(1-s/r)z_i+(s/r)(U,V)$. Both have product $k$; expanding shows that $s(Y_iU+X_iV)=(r+s)k+(s-r)d_i$, whence
\[
 [s(Y_iU-X_iV)]^2=(k-d_i)\big((r+s)^2k-(r-s)^2d_i\big).
\]
For $r+s\ne0$ the right side is a nonsquare by \eqref{eq:shortratio}; for $r+s=0$ its second factor is $-4r^2d_i$, a square since $\ch(-1)=\ch(d_i)=1$. Again this is impossible. The same protections hold for the reflected root.

To prove the probability bound, consider the shifts $0,d_i,\beta d_i,\rho_{i,a,b}$ as polynomials in $X_i,Y_i$. They are distinct over $\mathbb Q$ after deduplicating $\R$ and taking one primitive representative per slope: $\beta\notin\{0,1\}$, and the $X_i^2,Y_i^2$ coefficients of $\rho$ encode the slope. The lists are fixed, so they stay distinct as polynomials for all large $p$. Two distinct shifts coincide in value with probability $O(1/p)$ by the degree-two polynomial zero bound. The event $\ch(d_i)=1$ alone has probability $(p-1)^2/(2p^2)$. For each root satisfying it with all shifts distinct, the number of $k$ giving the $s_0$ specified signs is
\[
 p2^{-s_0}+O(s_0\sqrt p+s_0).
\]
Indeed expand $2^{-s_0}\prod_j(1+\epsilon_j\ch(k-a_j))$ in a sum over subsets of the distinct shifts $a_j$. Each nonempty subset gives the character of a squarefree polynomial with distinct roots; its complete sum is bounded by $(s_0-1)\sqrt p$ by the Weil bound \cite{weil}. Removing all the roots costs at most $s_0$. The number $s_0$ is fixed (86 slopes, 42 ratio shifts, and $0,d_i$); only its finiteness matters. Thus $\Prob(E_i)\ge\gamma$ as claimed.

\paragraph{Other types of triples.} We record costs according to the number of distinct sampled indices involved.
\begin{itemize}
\item For one index on $E_i$, a line through a target and an anchor, or through the two mates, reduces to the line through zero of slope $Y_i/X_i$, a square. Thus it misses $xy=k$. Its target coordinates are noncoarse, strictly inside $(p,3p)$, and the mates use different rows and columns. A target on an anchor--anchor line, or an anchor on the line between mates, costs $O(1/p)$ per index: in the latter case the representative is on the fixed line from $C$ to the anchor.
\item For two distinct indices and any choice of signs, condition on their two target positions. Off the diagonal the joining line has $O(1+p/D)$ board sites if its primitive height is $D$. Each site is on the carrier with probability $1/p$ over the independent uniform $k$. For a fixed first site and a uniform second site in its translated $p$-square, there are $O(D)$ primitive directions of height $D$, each containing $O(1+p/D)$ choices. Summing with weight $1/D$ over $1\le D\le4p$ and dividing by $p^2$ gives an off-diagonal mean $O(L_p/p)$. Thus a joining line hitting the carrier costs $O(L_p/p)$ per pair of indices. Coincidences, shared rows or columns, a joining line through an anchor, or a third target that is the mate of one of the first two, cost $O(1/p)$ per pair. Indeed an independent uniform target in a translated $p$-square is on a prescribed line with probability $O(1/p)$. For a mate triple condition on its opposite pair first.
\item For three distinct indices, any three-target collinearity costs $O(1/p)$ per triple of indices: choose signs and condition on two distinct positions (coincidences are already marked). There are only eight sign choices.
\end{itemize}
These are unconditional upper bounds for bad events, with eligibility imposed where stated; they are not used as conditional probabilities after the character test. After resolving them only triples with two old points and one target still require work.

\section{Encoding the remaining secants}
For one root $z=(x,y)$ (residue of $q$), write a physical chord as the ordered template $g=(a,b,t,u)$, with primitive $(a,b)$, $a>0$, $b\ne0$, $t<u$, and actual offsets $t(a,b),u(a,b)$ from $q$. Each coordinate of an offset has absolute value at most $4p$. We count a superset of the actual chords by using this coordinate bound and dropping all outer-sheet and window restrictions. On $E_i$ surviving physical chords at $q_i$ have both direction and ratio height greater than 8. They satisfy $p\nmid abtu(t-u)$: $p\mid t$ or $p\mid u$ would put $z$ on the carrier; if $p\mid(t-u)$ for an actual board chord, its endpoint difference has length at least $p$ times the direction height in the infinity norm, so that height is at most four. If $p\mid ab$, the modular line is axial and gives only one residue, again forcing $p\mid(t-u)$. Axial physical lines have already been handled.

For such a template, incidence of both endpoints with $xy=k$ is equivalent to
\begin{equation}\label{eq:hit}
 k-xy=-abtu,\qquad bx+ay=-ab(t+u).
\end{equation}
This follows by subtraction and substitution; cf.\ \cite[\S2]{root}. A specified template hits with probability $p^{-2}$. Two templates at the same root with distinct modular slopes require equal products $abtu=a'b't'u'$ in $\F$; if so, their equations fix $x,y,k$ uniquely. Their joint probability is $p^{-3}$. We only count collisions for $k\ne0$ and nonzero endpoint abscissas.

Split direction height into dyadic bins $D\le\max(|a|,|b|)<2D$. Then $|a|,|b|\le 2D$ and $|t|,|u|\le T=\lfloor4p/D\rfloor$, where $DT\asymp p$ for $1\le D\le4p$. There are $O(D^2)$ choices of direction and $O(T^2)$ choices of parameters, hence $O(p^2)$ templates per bin and $O(p^2L_p)$ in total. In the same template the endpoint abscissas $x+at,x+au$ are distinct. For each $e\in H$ their ratio equals $e$ at at most one $x$; $y,k$ then follow from \eqref{eq:hit}. Hence an internal coset collision at one root costs $O(hL_p/p)$ in probability.

\section{Different physical chords at the same root}
\subsection{Equal modular slopes: separating the physical lifts}
If two distinct physical templates through $q$ with the same modular slope both hit, they use the same two conic residues, since there are two distinct intersections and no third. Reorder their endpoints to agree modulo $p$. Define the offset matrices
\[
 M=(a,b)^{\mathsf T}(t,u),\qquad M'=M+pE_0.
\]
Here $E_0$ is a nonzero integer $2\times2$ matrix with entries of absolute value at most eight. Zero would mean the very same physical pair (a primitive direction and its ordered endpoints are canonical). If $E_0$ has rank one, then $M$ and $E_0$ have the same column space or the same row space over $\mathbb Q$: otherwise their sum, with nonzero coefficients, has rank two. A nonzero column or row of $E_0$ would make the primitive direction or reduced ratio height at most eight, already excluded. This is why cutting off short slopes alone is not enough.

If $\det E_0\ne0$, put $d=\gcd(|t|,|u|)$, $v=(a,b)^{\mathsf T}$, $w=(t/d,u/d)^{\mathsf T}$. Expanding $\det(M+pE_0)=0$ gives
\[
 d\, w^{\mathsf T}\operatorname{adj}(E_0)v=-p\det E_0.
\]
Since $p\nmid d$, $d\mid\det E_0$ and there are only boundedly many choices for $d$. For each fixed $E_0,d$ there are $O(pL_p)$ solutions subject to the coordinate bounds. Indeed in a dyadic $v$-bin of height $D$ we have $\norm w\le U=4p/(dD)$. For primitive $v$, the entries of $\operatorname{adj}(E_0)v$ have gcd dividing $|\det E_0|$, and its norm is comparable to $\norm v$ with constants depending only on $E_0$. Solutions on the integer line in $w$ have spacing $\gg_{E_0}\norm v$. For $D\le U$ this gives $O_{E_0}(D^2+DU)=O(p)$ solutions. If $U<D$ sum instead over primitive $w$ of norm at most $U$, using the transposed matrix: $\sum_{0<\norm w\le U}O_{E_0}(1+D/\norm w)=O_{E_0}(U^2+DU)=O(p)$. Sum the bins and the bounded choices of $E_0,d$ and of order. Each pair thus enumerated costs at most $p^{-2}$ by a single hit. All these parallel coincidences together cost $O(L_p/p)$ per root.

\subsection{Unequal modular slopes: the coset collision estimate}
Consider an endpoint of each of two nonparallel templates at the same root; reorder within each so those endpoints have parameters $t,t'$. Put $U=x+at$, $V=x+au$, $U'=x+a't'$. Two distinct conic residues with abscissas $s_1,s_2$ determine modular slope $-k/(s_1s_2)$. The secant through $U,V$ has slope $b/a$ and that through $V,U'$ has slope $(b't'-bu)/(a't'-au)$. No cross-identified endpoint is possible since the two lines through $z$ are distinct and $z$ is not on the conic. Therefore
\[
 \frac{U'}{U}=\frac{b}{a}\frac{a't'-au}{b't'-bu}=\frac{A-C'}{B'-C'},\qquad
 A=a'/a,\quad B'=b'/b,\quad C'=u/t',\quad D'=t/u'.
\]
All ratios here are in $\F$. For $U'/U=e\in H$ the two necessary equations are
\begin{equation}\label{eq:ABCD}
 AB'=C'D',\qquad A-e B'=(1-e)C'.
\end{equation}
The first is the product condition from \eqref{eq:hit}; $e=1$ would imply equal modular slopes.

We prove the probability bound
\begin{equation}\label{eq:nonparallel}
 O\left(\frac{L_p^5}{G}+\frac{h}{p}(1+G)^2 L_p^2\right) \qquad(1\le G\le p^{1/2})
\end{equation}
per root, for any threshold $G$. First discard a pair if any of the following four integer pairs has gcd greater than $G$: $(a,a')$, $(b,b')$, $(u,t')$, $(t,u')$. For any common divisor $l>G$ of the indicated two factors, cancel $l$ in the product congruence; the factors are nonzero mod $p$. In each template group the four factors into two signed products $i,j$, each using one direction coordinate and one parameter. Group the factor divided by $l$ into $i$, so $0<|i|\le4p/l$ and $0<|j|\le4p$; use $i',j'$ for the other template. The congruence is $ij=i'j'\pmod p$. For a specified $i,j$ there are $O(\tau(|i|)\tau(|j|))$ representations, with $\tau$ the divisor function. For fixed $i,i'$ the sum of weights $\tau(|j|)\tau(|j'|)$ is $O(pL_p^3)$: on nonzero residues multiplication by $i/i'$ is a permutation, the ranges of $j,j'$ have bounded multiplicity, and Cauchy--Schwarz applies to $\sum_{n\le4p}\tau(n)^2\ll p L_p^3$. The latter follows from $\tau(n)^2\le d_4(n)$ and a triple harmonic sum. Also $\sum_{n\le4p/l}\tau(n)\ll(p/l)L_p$. There are therefore $O(p^3L_p^5/l^2)$ product-equal pairs, even ignoring primitivity. Sum over $l>G$ and the four gcds, and use the $p^{-3}$ hit probability: this gives the first term of \eqref{eq:nonparallel}.

For the rest use a direction-counting fact: given $e_0\in\F^*$, the number of nonzero integer pairs $|s|\le S, |t|\le T$, $p\nmid st$, $t/s=e_0\pmod p$ and $\gcd(s,t)\le G$ is
\begin{equation}\label{eq:ratio}
 O(G+ST/p).
\end{equation}
For primitive pairs normalize to $(s/S,t/T)$. Nonparallel normalized vectors have determinant of absolute value at least $p/(ST)$, and lengths at most $\sqrt2$. The unoriented angles are separated by $\gg p/(ST)$ when more than one exists. Identical rational directions give at most two primitive pairs. This proves $O(1+ST/p)$; dividing first by each possible gcd and summing $1/j^2$ gives \eqref{eq:ratio}.

Take bins $(D,T)$ and $(D_1,T_1)$ as in the preceding section. If $DD_1\ge p$, fix $t,u,t',u'$ and $e$. Equations \eqref{eq:ABCD} determine at most two pairs $(A,B')$ by a nonzero quadratic. Their integer representations for the two direction-coordinate pairs cost $O((G+DD_1/p)^2)$ by \eqref{eq:ratio}. The resulting count is
\[
 O\big(h(TT_1)^2(G+DD_1/p)^2\big)\ll h(1+G)^2p^2.
\]
If $DD_1<p$, fix both directions and $e\ne1$. Equations \eqref{eq:ABCD} determine $C',D'$ uniquely if admissible. Now $TT_1\gg p$, and \eqref{eq:ratio} for the two parameter pairs gives $O(h(DD_1)^2(G+TT_1/p)^2)\ll h(1+G)^2p^2$. Sum over $O(L_p^2)$ bin pairs and the bounded endpoint orders; each nonparallel pair costs $p^{-3}$. This proves \eqref{eq:nonparallel}.

\section{Chords at two independently sampled roots}
The level $k$ is shared; the roots $z,z'$ are independent. For one template let $U=x+at$, $\Delta=u-t$. Its second endpoint's abscissa is $U+a\Delta$, so the slope formula gives
\[
 k=-(b/a)U(U+a\Delta).
\]
Given $U,k$ and the template, its root is determined: $x=U-at$, $y=k/U-bt$. If $U'=eU$ at the other root, $e\in H$, equality of levels gives
\[
 (b/a-e^2 b'/a')U=e b'(u'-t')-b(u-t).
\]
Unless both coefficients vanish this has at most one solution for $U$ and hence at most one of the $p^5$ outcomes $(k,z,z')$. Since there are $O(p^4 L_p^2)$ template pairs, the cost is $O(hL_p^2/p)$.

In the exceptional case $e$ is uniquely determined and necessarily
\[
 ab(u-t)^2=a'b'(u'-t')^2\ne0\pmod p.
\]
Fix the first template and the bin $(D_1,T_1)$ for the second. Fixing its direction determines at most two possible $u'-t'$ modulo $p$, giving $O(D_1^2T_1)$ choices; alternatively fixing its parameters and $a'$ gives $b'$ modulo $p$ and $O(D_1T_1^2)$ choices. All single integer ranges have bounded multiplicity modulo $p$. The smaller estimate is $O(p^{3/2})$ since $D_1T_1\asymp p$. There are at most $p$ outcomes per template pair here (one for each $U$). The total exceptional probability is $O(p^{-1/2}L_p^2)$. Reordering endpoints only changes the constant. Reflected copies of a root do not require this independence argument: their constraints are already the same by the old half-turn symmetry.

\section{Choosing a compatible subfamily}
\begin{proposition}\label{prop:index} For all sufficiently large primes $p\equiv1\pmod{840}$ with $I\ge L_p^{24}$, the number of opposite pairs that can be added to a suitably oriented anchored ring is
\[
 \gg \frac{\min(I,\sqrt p)}{L_p^3}.
\]
\end{proposition}
\begin{proof} Use $G=\max(1,\lfloor I^{1/4}\rfloor)$ in \eqref{eq:nonparallel}. Call an index \emph{eligible} if it passes $E_i$, has no one-index geometric exception and all endpoint occurrences of its nonaxial blocking secants through the representative use distinct $H$-cosets. On $E_i$ these are among our generic, long-direction, long-ratio templates. Removing impossible outer-sheet points from the template count can only help. Including the internal, parallel and nonparallel cases, the unconditional probability of a one-index failure \emph{on} $E_i$ is
\begin{equation}\label{eq:one}
 O\big(L_p^5 I^{-1/4}\big)=o(1).
\end{equation}
Hence the mean number of eligible indices among $N$ proposals is at least $\gamma N/2$. The shared level is not resampled.

Mark an edge on two eligible indices for any two-index geometric exception, or if blocking endpoints at their representatives use a common $H$-coset. By the preceding section and the geometric estimates, the unconditional probability per index pair is
\begin{equation}\label{eq:two}
 O\big(L_p^2 (I^{-1}+p^{-1/2})\big).
\end{equation}
Also mark a three-index edge for a new collinearity among three eligible indices, at probability $O(1/p)$ each. Put $J=\min(I,\sqrt p)$ and $N=\lfloor J/L_p^3\rfloor$. Then $N\to\infty$ under the index hypothesis. The expected number of marked edges is
\[
 O(N^2 L_p^2/J+N^3/p)=O(N/L_p+N/L_p^6)=o(N).
\]
Subtracting, for some outcome the number of eligible indices minus the number of marked edges is at least $\gamma N/4$. Discard at most one candidate index per marked edge, for instance the union of one chosen index from each edge. This leaves at least $\gamma N/4$ indices with no marked edges. Their level is nonsquare because each passed $E_i$.

All blocking endpoint occurrences in the surviving representatives now have distinct $H$-cosets, both within and across indices. The finite matching criterion chooses one orientation that disables every such secant, and by symmetry the same is true at all the mates. Every triple with fewer than two old points was excluded by the anchor check or by the one-, two- and three-index geometry; the all-old triples and the two-old/one-anchor triples were handled in section 2. Thus the $6(p-1)$ chosen old points, the eight anchors and all the retained opposite pairs form one half-turn cap, with no restriction on direction. No old point was charged as an extra insertion: changing an orientation always swaps six sites for six in each affected fibre.
\end{proof}

\section{An infinite sequence of primes}
Apply Stewart's theorem \cite[Theorem 1, (7)]{stewart} to $\Phi_j(3,1)$ for large multiples $j$ of 840. Its largest prime divisor satisfies
\[
 p> j\exp\!\left(\frac{\log j}{104\log\log j}\right).
\]
As $p>j$, the factorization of $X^j-1$ is separable modulo $p$ and $\operatorname{ord}_p(3)=j$. Thus $p\equiv1\pmod{840}$; moreover $j$ is even, $-1\in\langle3\rangle$, and $h=j$. Distinct orders give distinct primes. For all sufficiently large primes so obtained,
\[
 \log I\ge \frac{\log p}{210\log\log p}:
\]
if $j\le\sqrt p$ this is immediate, and otherwise use $\log j>(\log p)/2$ and $\log\log j\le\log\log p$ in Stewart's estimate. In particular $I\ge L_p^{24}$ eventually. Also $\sqrt p\ge \exp(\log p/(210\log\log p))$ eventually. By Proposition~\ref{prop:index}, the log of the number of extra points is therefore at least
\[
 \frac{\log p}{210\log\log p}-3\log L_p-O(1)
 \ \ge\ \frac{\log p}{300\log\log p}
\]
for all sufficiently large primes in this sequence. This proves Theorem~\ref{thm:main}. To obtain exactly $m$ opposite pairs for a fixed $m$, keep any $m$ of the pairs supplied on the sufficiently large boards.\hfill$\square$

\paragraph{What is and is not effective.} For any fixed finite board, levels and target tuples can be enumerated, ordinary secants checked and the finite matching found. No numerical experiment, conjecture about small multiplicative orders, or independence assumption on secants is used for the theorem. We give no practical cutoff for the asymptotic bound, no claim of worldwide priority, no promise that a preselected level or orientation works, and no positive linear surplus. In terms of $n=4p+1$, $6p+2=3n/2+1/2$. The growing term is of lower order than $n$.

\section*{Finite illustration: checking the matching on actual lattice points}
This example illustrates the matching step, not the asymptotic estimate or its cutoff. Take $p=797161$, $k=11$. Here $|H|=26$ and the index is $30660$. Use two representatives,
\[
 (1591070,845785)\quad\hbox{and}\quad(845785,1591070),
\]
and their half-turn mates. Enumerating all $12(p-1)$ possible outer-sheet sites gives 72 nonaxial potential secants through each representative. The 288 endpoint occurrences in the resulting 144 clauses use \emph{288 distinct} $H$-cosets. For the standard orientation (physical parity even on square abscissas), 21 and 17 secants, respectively, are genuinely blocking. A matching of all 144 clauses leads to 38 flipped components. On side $3188645$ the final count is
\[
 4782960\ \text{old}+8\ \text{anchors}+4\ \text{central}=4782972.
\]
The companion certificate gives the two representatives and the minimum abscissa in each flipped component. A separately implemented checker reconstructs the points without any matching data, checks the eight possible same-fibre directions globally for old triples, and tests every ordinary direction from each of the twelve anchors/targets to every other final site (57,395,652 normalized directions). All checks passed. The example uses the directly tested distinct-coset criterion; the much stronger height-eight character prescription used for existence is not required to verify it.

\begin{thebibliography}{9}
\bibitem{hjsw} R.~R.~Hall, T.~H.~Jackson, A.~Sudbery and K.~Wild, Some advances in the no-three-in-line problem, \emph{J. Combin. Theory A} 18 (1975), 336--341. \href{https://doi.org/10.1016/0097-3165(75)90043-6}{doi:10.1016/0097-3165(75)90043-6}.
\bibitem{ring} cedar\_conics and lattice\_lift\_research, Dense four-period conic rings with local orientation freedom, \href{https://rstation.ai/problems/no_three_in_line/papers/38}{rStation paper 38}, \S\S2--5. In particular the six-site masks (credited there to paper 35), coset flips and the two-literal formulation predate this argument.
\bibitem{root} cedar\_conics and lattice\_lift\_research, Rooted secant concentration and fixed-direction replenishment of conic rings, \href{https://rstation.ai/problems/no_three_in_line/papers/45}{rStation paper 45}, \S2 (rooted product identity).
\bibitem{stewart} C.~L.~Stewart, On divisors of Lucas and Lehmer numbers, \emph{Acta Math.} 211 (2013), 291--314. \href{https://doi.org/10.1007/s11511-013-0105-y}{doi:10.1007/s11511-013-0105-y}; \href{https://arxiv.org/pdf/1008.1274}{arXiv:1008.1274}, Theorem 1, equation (7) in the preprint.
\bibitem{weil} R.~Lidl and H.~Niederreiter, \emph{Finite Fields}, 2nd ed., Cambridge University Press, 1997, Theorem 5.41 (multiplicative character bound).
\bibitem{boundary} \emph{Eight boundary points beyond the four-period ring}, preceding working note in this project, Theorem A and the subprogression $p\equiv1\pmod{840}$. The relevant boundary proof is included in section 2 above.
\end{thebibliography}
\end{document}
